\chapter{Discussion, Assumptions, and Limitations}
\label{ch:discussion}

\section{Explicit Modeling Assumptions}
\label{sec:disc-assumptions}

The assignment specification leaves a small number of quantities
underspecified. Every place this report's implementation makes a concrete
choice is collected here.

\begin{itemize}
    \item \textbf{$T_{\mathrm{ref}}$ in the cost function} (\cref{eq:cost}):
          taken as the tightest configured required time among the
          circuit's primary outputs, rather than, for example, a single
          global clock period constant.
    \item \textbf{Tie-breaking in candidate ranking}
          (\cref{sec:opt-swc,sec:opt-ceg}): ties in a heuristic's score are
          broken by netlist declaration order, for reproducibility, not by
          any physical criterion.
    \item \textbf{Short-circuit power}: not modeled, per the specification's
          explicit permission to simplify (\cref{sec:bg-power}).
    \item \textbf{Non-unate required-time propagation}
          (\cref{sec:sta-backward}): both output transitions are treated as
          simultaneously legal for a non-unate pin, and the more
          restrictive (minimum) required time is applied to both the rise
          and fall required time propagated backward through that pin --- a
          conservative reading of the specification's non-unate timing-arc
          description, chosen because it cannot understate a required-time
          constraint.
\end{itemize}

\section{Known Limitations}
\label{sec:disc-limitations}

\begin{itemize}
    \item \textbf{Brute force is not the canonical pipeline heuristic.}
          It is the most thorough of the four (\cref{sec:opt-bruteforce})
          but is only invoked manually or by the test suite
          (\cref{sec:test-equivalence}), not by the standard
          \code{app/experiments.py} pipeline that produces the results in
          \cref{ch:results}. A reader comparing ``the tool's result'' to
          brute force's own result should be aware these are not the same
          run.
    \item \textbf{The beam-search benchmark reference is not certified
          globally optimal} (\cref{sec:bench-reference}): it is explicitly
          labeled \code{best\_known\_multi\_start\_beam} in the generated
          artifacts, and evaluation metrics that compare a heuristic
          against it should be read as ``distance from the best solution
          found by a substantially less restricted search,'' not ``distance
          from a proven optimum.''
    \item \textbf{Single clock domain / purely combinational assumption.}
          Consistent with the assignment scope, no sequential elements,
          multiple clock domains, or setup/hold-time analysis at a
          register boundary are modeled; every primary input and output
          required/arrival time is treated as an independent, externally
          supplied constraint.
    \item \textbf{Sequential, first-improvement search order in three of
          the four heuristics} (\cref{sec:opt-heuristics}) means
          slack-weighted capacitance, criticality/effort-gap, and random
          greedy can all, in principle, accept an earlier change that
          later turns out to block a more valuable one on the area/power
          budget --- \cref{sec:res-case-divergence} documents exactly this
          happening to criticality/effort-gap on \code{c13}. Brute force's
          best-of-batch search does not have this failure mode, at the
          cost of evaluating substantially more candidates per iteration.
\end{itemize}

\section{Threats to Validity for the Benchmark-Suite Results}
\label{sec:disc-threats}

\begin{authornote}
Expand this section once \cref{sec:res-benchmarking} has been completed
with actual synthetic-suite results. Points worth addressing explicitly:
whether the procedural DAG generator's structural distribution (size,
depth, fan-out) is representative of realistic combinational logic or
biased toward easier/harder-than-typical cases; whether the planted-
downsizing mutation procedure (\cref{sec:bench-planted}) produces
violations with a severity distribution comparable to the fixed benchmark
suite's; and whether Random Greedy's repeated-seed evaluation
(\cref{sec:bench-parallel}) uses enough seeds per case for its reported
metrics to be statistically stable, given the observed variance.
\end{authornote}

\section{Threats to Validity for the Fixed-Suite Comparison}
\label{sec:disc-fixed-threats}

The fixed-suite comparison in \cref{sec:res-heuristics} is based on fifteen
circuits, six of which exhibit a violation; this is, by design, the
motivation for \cref{ch:benchmarking}'s synthetic suite rather than a
statistically independent confirmation of it. The two case studies in
\cref{sec:res-case-c02,sec:res-case-divergence} are illustrative of a
specific, verified mechanism (ordering interacting with the area/power
budget) rather than a claim that this mechanism explains every observed
difference between heuristics across the whole suite; a reader should treat
them as worked examples of \emph{how} the heuristics can diverge, not as an
exhaustive causal account of \cref{tab:res-heuristics} as a whole.
